% ch10_a 第10章 §10.1–§10.2 (书页 440–452 / p455–p467)
% 第10章 STRING CONDENSATION—AN UNIFICATION OF LIGHT AND FERMIONS

\chapter{弦凝聚——光与费米子的统一}

\begin{keypoints}
\item \textbf{光与费米子是什么？} 光是任意大小的、由凝聚弦构成的网的涨落。费米子是凝聚弦的端点。
\item \textbf{光与费米子从何而来？} 光与费米子来自充满空间的凝聚弦网的集体运动。
\item \textbf{光与费米子为何存在？} 光与费米子之所以存在，是因为我们的真空恰好具有弦网凝聚（string-net condensation）。
\end{keypoints}

在第 9 章中，我们用投影构造在二维自旋系统中构造了许多量子有序态，并引入投影对称群（PSG）来表征和分类这些量子有序态。这些量子有序态不仅包含规范玻色子，而且还包含作为其低能集体激发的费米子。规范玻色子和费米子的演生是一个非常惊人的结果，因为底层的格子模型是一个纯粹的玻色模型。

许多年来，费米子和规范玻色子一直被视为基本的、不可触碰的。当有人提出费米子和规范理论可以从二维自旋液体中演生出来时（Baskaran et al., 1987; Kivelson et al., 1987; Baskaran and Anderson, 1988），这一提议甚至没有引起惊讶，因为它遭遇到的是怀疑与不信。这些怀疑是有根据的。即使到了现在，我们仍不清楚二维自旋液体是否存在。然而，这并不意味着最初的提议是错误的。把这种计算应用到 $SU(N)$ 自旋模型上时（Affleck and Marston, 1988），它的确在三维中给出良好定义的规范玻色子和费米子（见 10.7 节）（Wen, 2002a）。因此，费米子和规范玻色子并非那样基本而不可触碰。它们可以从卑微的玻色模型中的某些量子序演生出来。

然而，导致规范玻色子和费米子演生的投影构造是非常形式化的。似乎一切都依赖于一个非常不合情理的数学技巧——把一个自旋拆成两个非物理的费米子。人们很难对这样一种方法抱有任何信心。由这种方法得到的演生费米子和规范玻色子这一惊人结果，看起来实在难以置信。

本章可以被看作一个建立信心的练习。我们将讨论几个精确可解的自旋模型，它们同样可以用投影构造来求解。我们将证明，对那些模型，投影构造给出的结果与精确结果一致。特别地，正如投影构造所提示的，这些精确可解模型具有演生的费米子和一个规范场。

然而，更重要的是，精确可解模型把费米子和规范场的演生与一种现象——闭合弦网的凝聚——联系在了一起。我们发现，由投影构造得到的量子有序态实际上就是弦网凝聚态。表征由投影构造得到的态中的量子序的 PSG，实际上表征的是不同的弦网凝聚。投影构造与弦网凝聚只是描述同一类量子序的两种不同方式。弦网凝聚为形式化的投影构造提供了一个物理基础。

试图联系投影构造与弦网凝聚，只不过是本章诸讨论的一个形式上的动机。既然费米子和规范玻色子可以从弦网凝聚中演生，我们想为本章的讨论赋予一个更物理的动机。我们想问：弦网凝聚与我们宇宙中光和费米子的存在有没有关系？特别是，我们想就光与费米子提出下面这些问题。光与费米子是什么？光与费米子从何而来？光与费米子为何存在？

目前，对上述基本问题的标准回答似乎是“光是由规范场描写的粒子”以及“费米子是由反对易场描写的粒子”。在这里，我们想论证，对上述问题还存在另一种可能的（而且我相信是更物理的）回答，即：我们的真空充满了任意大小的、由弦状物体构成的网，而这些弦网形成一个量子凝聚态。按照弦网理论，光（以及其他规范玻色子）是凝聚弦网的集体振动，而费米子是弦的端点。我们看到，弦网凝聚为光和费米子提供了一个统一的起源。\footnote{这里，“弦网凝聚”指的是由任意大小的弦状物体构成的网的凝聚。}换句话说，弦网凝聚统一了光与费米子。如果有人说“要有弦网”，那么我们就同时得到了光和费米子。

在为光与费米子的弦网理论提供证据之前，我们想先澄清所谓“光存在”与“费米子存在”是什么意思。我们知道，物理中有一个天然的质量标度——普朗克质量。普朗克质量非常之大，以致我们看到的任何粒子，其质量至少比普朗克质量小 $10^{16}$ 倍。因此，与普朗克质量相比，所有观测到的粒子都可以当作无质量的来处理。当我们问某些粒子为什么存在时，我们的真正意思是：与普朗克质量相比，那些粒子为什么是无质量的（或近乎无质量的）。真正的问题是什么使某些类型的激发（例如光和费米子）成为无质量的（或近乎无质量的）。自然界究竟为什么想要拥有无质量的激发？是谁订购了它们？

其次，我们想澄清所谓“光与费米子的起源”是什么意思。我们知道，一切都必须来自某种东西。所以当我们问“光与费米子从何而来？”时，我们已经假定了存在比光和费米子更简单、更基本的东西。在 10.1 节中，我们将定义局域玻色模型（local bosonic model），它比具有规范场与费米子耦合的模型更简单。我们将把局域玻色模型看作更基本的。这样一种哲学将被称为局域性原理（locality principle）。我们将证明，如果一个局域玻色模型的基态中含有由弦状物体构成的网的凝聚，那么光和费米子就可以从这个模型中演生出来。

费米子的弦网理论解释了为什么我们的宇宙中费米子总是成偶数地出现——因为一条弦（或一个弦网）总有偶数个端点。关于规范玻色子和费米子的弦网理论还有一个实验预言，即所有费米子都必须携带某种规范电荷。乍看之下，这个预言似乎与已知的实验事实——中子不携带任何规范电荷——相矛盾。于是有人可能认为，关于规范玻色子和费米子的弦网理论已经被实验证伪了。这里我们想指出，如果我们假定我们的宇宙中还存在一种新的分立规范场，例如一个 $Z_2$ 规范场，那么关于规范玻色子和费米子的弦网理论仍然可以是正确的。在这种情形下，中子和中微子携带该分立规范场的一个非零电荷。因此，关于规范玻色子和费米子的弦网理论预言，我们的宇宙中存在分立的规范激发（例如规范通量线）。

按照量子序与弦网凝聚的图像，基本粒子（例如光子和电子）可能终究不是基本的。它们可能是普朗克标度以下某个局域玻色系统的集体激发。由于我们无法在接近普朗克标度处做实验，很难判定光子和电子究竟是不是基本粒子。在本章中，我们将通过二维和三维格子上的一些具体的局域玻色模型，证明光与费米子的弦网理论至少是自洽的。这里研究的局域玻色模型，只是长长的局域玻色模型清单中的几个例子，这些模型都包含演生的非禁闭费米子和规范玻色子（Foerster et al., 1980; Kalmeyer and Laughlin, 1987; Wen et al., 1989; Read and Sachdev, 1991; Wen, 1991a; Moessner and Sondhi, 2001; Balents et al., 2002; Ioffe et al., 2002; Motrunich and Senthil, 2002; Sachdev and Park, 2002; Kitaev, 2003; Motrunich, 2003; Wen, 2003c）。这些模型的基态都具有非平凡的拓扑序/量子序。

\begin{smallnote}
我们想指出，尽管有某些相似之处，上述关于规范玻色子和费米子的弦网理论不同于标准的超弦理论。在标准超弦理论中，闭合弦对应于引力子，开弦对应于规范玻色子。费米子来自世界面上的费米型场。超弦理论中所有基本粒子（包括规范玻色子）都对应于小弦的不同振动模式。而在弦网理论中，真空充满了大弦（或大弦构成的网，见图 1.2）。无质量规范玻色子对应于大闭合弦之网的涨落，费米子对应于开弦的端点。弦网理论中没有费米型场。

关于规范涨落的弦网理论与规范理论中的 Wegner--Wilson 圈密切相关（Wegner, 1971; Wilson, 1974; Kogut, 1979）。动力学规范理论与动力学 Wegner--Wilson 圈理论之间的关系由 Gliozzi et al.~(1979) 和 Polyakov (1979) 提出并研究。Savit (1980) 综述了规范理论与延展物体理论之间的各种对偶关系。特别地，人们发现某些统计格点规范模型对偶于某些统计膜模型（Banks et al., 1977）。这种对偶关系与量子模型中规范理论和弦网理论之间的关系直接相连。量子模型中的新特点是，弦的端点有时可以是费米子或任意子（Levin and Wen, 2003）。
\end{smallnote}

这里我们想强调，对于我们宇宙中实际的规范玻色子和费米子而言，弦网理论目前只是一个建议。尽管弦网凝聚能够产生并保护无质量的光子、胶子、夸克以及其他带电轻子，但我们目前还不知道弦网凝聚能否产生中微子——一种手征费米子。我们也不知道弦网凝聚能否产生弱相互作用的 $SU(2)$ 规范场——它与夸克和轻子手征地耦合。弦网凝聚在我们真空中的正确性，取决于上述问题的解决。

另一方面，如果我们只关心凝聚态物理的问题——如何用玻色子制造人工光和人工费米子——那么弦网理论和量子序确实给出了一个答案。要制造人工光和人工费米子，我们只需让某些弦凝聚。

\section{局域玻色模型}

在本章中，我们将只考虑局域玻色模型。我们认为局域玻色模型是基本的，因为它们是真正局域的。我们注意到，费米子模型一般是非局域的，因为不同格点上的费米子算符即使在格点相距很远时也不对易。与此相反，局域玻色模型是局域的，因为相距很远的玻色子算符彼此对易。由于局域玻色模型具有内在的局域性，我们相信自然的基本理论是一个局域玻色模型。为强调这一点，我们将给这一信念起一个名字：局域性原理。

让我们详细定义局域玻色模型。为了定义一个物理系统，我们需要明确（i）一个总希尔伯特空间，（ii）一组局域物理算符的定义，以及（iii）一个哈密顿量。基于这一理解，局域玻色模型定义为满足下列性质的模型：

\begin{enumerate}
\item[(i)] 总希尔伯特空间是有限维局域希尔伯特空间的直积。
\item[(ii)] 局域物理算符是作用在单个局域希尔伯特空间之内的算符，或是这些算符对邻近局域希尔伯特空间的有限乘积。我们也将称这些算符为局域玻色算符，因为相距很远时它们彼此对易。
\item[(iii)] 哈密顿量是局域物理算符之和。
\end{enumerate}

格子上的自旋 $1/2$ 系统是局域玻色模型的一个例子。其局域希尔伯特空间是二维的，包含 $|\uparrow\rangle$ 和 $|\downarrow\rangle$ 态。局域物理算符是 $\sigma_{\pmb{i}}^{a}$、$\sigma_{\pmb{i}}^{a}\sigma_{\pmb{i}+\pmb{x}}^{b}$ 等，其中 $\sigma^{a}$（$a=1,2,3$）是泡利矩阵。

自由的无自旋费米子系统（二维或更高维）不是局域玻色模型，尽管它拥有与自旋 $1/2$ 系统相同的总希尔伯特空间。这是因为不同格点上的费米子算符 $c_i$ 不对易，因而不是局域玻色算符。更重要的是，费米子跳跃算符 $c_i^\dagger c_j$ 在二维及更高维中不能写成局域玻色算符。（不过，由于 Jordan--Wigner 变换，一维费米子跳跃算符 $c_{i+1}^\dagger c_i$ 可以写成局域玻色算符。因此，如果把 $c_i$ 排除在局域物理算符的定义之外，一维费米系统可以是局域玻色模型。）格点规范理论不是局域玻色模型，因为它的总希尔伯特空间不能是局域希尔伯特空间的直积。

\section{由投影构造得到的一个精确可解模型}

在本节中，我们将构造一个精确可解的局域玻色模型。在下一节中，我们将证明该模型包含弦网凝聚和演生的费米子。

\subsection{精确可解模型的构造}

\begin{keypoints}
\item 精确可解模型可以通过寻找彼此对易的算符来构造。
\end{keypoints}

通常，投影构造并不给出精确结果。在本节中，我们将在二维正方格子上构造一个精确可解模型（Kitaev, 2003; Wen, 2003c）。我们的模型具有如下性质：投影构造给出精确的基态以及所有其他精确的激发态。构造的关键一步是找到一组彼此对易的算符。

让我们引入四个马约拉纳费米子（Majorana fermion）算符 $\lambda_{\pmb{i}}^{a}$（$a=x,\bar{x},y,\bar{y}$），满足 $\{\lambda_{a,\pmb{i}},\lambda_{b,\pmb{j}}\}=2\delta_{ab}\delta_{\pmb{i}\pmb{j}}$。我们发现，算符
\begin{equation*}
\hat U_{\pmb{i},\pmb{i}+\hat{\pmb{x}}}=\lambda_{\pmb{i}}^{x}\lambda_{\pmb{i}+\hat{\pmb{x}}}^{\bar{x}},\qquad
\hat U_{\pmb{i},\pmb{i}+\hat{\pmb{y}}}=\lambda_{\pmb{i}}^{y}\lambda_{\pmb{i}+\hat{\pmb{y}}}^{\bar{y}},\qquad
\hat U_{\pmb{i}\pmb{j}}=-\hat U_{\pmb{j}\pmb{i}},
\tag{10.2.1}
\end{equation*}
构成一组彼此对易的算符，即 $[\hat U_{\pmb{i}_1\pmb{i}_2},\hat U_{\pmb{j}_1\pmb{j}_2}]=0$。在得到一组对易算符之后，我们容易看出，作为诸 $\hat U_{\pmb{i}\pmb{j}}$ 的函数，相互作用费米子哈密顿量
\begin{equation*}
H=-\sum_{\pmb{i}}V_{\pmb{i}}\hat F_{\pmb{i}},\qquad
\hat F_{\pmb{i}}=\hat U_{\pmb{i},\pmb{i}_1}\hat U_{\pmb{i}_1,\pmb{i}_2}\hat U_{\pmb{i}_2,\pmb{i}_3}\hat U_{\pmb{i}_3,\pmb{i}},
\tag{10.2.2}
\end{equation*}
与所有 $\hat U_{\pmb{i}\pmb{j}}$ 对易。这里 $\pmb{i}_1=\pmb{i}+\hat{\pmb{x}}$，$\pmb{i}_2=\pmb{i}+\hat{\pmb{x}}+\hat{\pmb{y}}$，而 $\pmb{i}_3=\pmb{i}+\hat{\pmb{y}}$。我们将把 $\hat F_{\pmb{i}}$ 称为 $Z_2$ 通量算符。

为了得到式 (10.2.2) 中的哈密顿量 $H$ 所作用的希尔伯特空间，我们在每个格点上把 $\lambda^{x,\bar{x},y,\bar{y}}$ 组合成两个复费米子算符
\begin{equation*}
2\psi_{1,\pmb{i}}=\lambda_{\pmb{i}}^{x}+\mathrm{i}\lambda_{\pmb{i}}^{\bar{x}},\qquad
2\psi_{2,\pmb{i}}=\lambda_{\pmb{i}}^{y}+\mathrm{i}\lambda_{\pmb{i}}^{\bar{y}}
\tag{10.2.3}
\end{equation*}
可以验证，$\psi_{1,2}$ 满足标准的费米子反对易代数。这两个复费米子算符在每个格点上生成一个四维希尔伯特空间。总希尔伯特空间的维数为 $4^{N_{\text{site}}}$。诸 $\hat U_{\pmb{i}\pmb{j}}$ 与 $H$ 的对易性使我们能够精确地求解这个相互作用费米子系统。

为了得到 $H$ 的精确本征态和精确本征值，令 $|\{s_{ij}\}\rangle$ 为诸 $\hat U_{ij}$ 算符的具有本征值 $s_{ij}$ 的共同本征态。由于 $(\hat U_{ij})^{2}=-1$ 且 $\hat U_{ij}=-\hat U_{ji}$，我们发现 $s_{ij}$ 满足 $s_{ij}=\pm\mathrm{i}$ 且 $s_{ij}=-s_{ji}$。由于 $H$ 是诸 $\hat U_{ij}$ 的函数，我们发现 $|\{s_{ij}\}\rangle$ 也是式 (10.2.2) 的一个能量本征态，其能量为
\begin{equation*}
E=-\sum_{\pmb{i}}V_{\pmb{i}}F_{\pmb{i}},\qquad
F_{\pmb{i}}=s_{\pmb{i},\pmb{i}_1}s_{\pmb{i}_1,\pmb{i}_2}s_{\pmb{i}_2,\pmb{i}_3}s_{\pmb{i}_3,\pmb{i}}
\tag{10.2.4}
\end{equation*}
我们将把 $F_{\pmb{i}}$ 称为穿过方格 $\pmb{i}$ 的 $Z_2$ 通量。当 $V_{\pmb{i}}>0$ 时，基态由 $|\{s_{ij}\}\rangle$ 给出，其中 $s_{\pmb{i},\pmb{i}+\pmb{x}}=s_{\pmb{i},\pmb{i}+\pmb{y}}=\mathrm{i}$。对这样的态，我们有 $F_{\pmb{i}}=1$。为了得到激发态，我们翻转某些 $s_{ij}$ 的符号，使某些 $F_{\pmb{i}}=-1$。

为了考察诸 $|\{s_{ij}\}\rangle$ 是否代表 $H$ 的所有精确本征态，我们需要清点态的数目。设二维正方格子有 $N_{\text{site}}$ 个格点，且两个方向都取周期性边界条件。这时格子有 $2N_{\text{site}}$ 条链。由于 $s_{ij}$ 总共有 $2^{2N_{\text{site}}}$ 种不同的选择（每条链有两种选择），态 $|\{s_{ij}\}\rangle$ 遍尽了希尔伯特空间中全部 $4^{N_{\text{site}}}$ 个态。因此，诸 $\hat U_{\pmb{i}\pmb{j}}$ 的共同本征态不简并，上述方法使我们能够得到 $H$ 的所有本征态和本征值。我们已经精确求解了这个二维相互作用费米子系统！

我们注意到，哈密顿量 $H$ 只能把每个格点上的费米子数改变偶数。因此，$H$ 作用在一个每个格点具有偶数个费米子的子空间之内。我们将称这个子空间为物理希尔伯特空间。物理希尔伯特空间在每个格点上只有两个态。限制在物理空间之内时，$H$ 实际上描写一个自旋 $1/2$ 或硬核玻色子系统。为了得到相应的自旋 $1/2$ 哈密顿量，我们注意到
\begin{equation*}
\sigma_{\pmb{i}}^{x}=\mathrm{i}\lambda_{\pmb{i}}^{y}\lambda_{\pmb{i}}^{x},\qquad
\sigma_{\pmb{i}}^{y}=\mathrm{i}\lambda_{\pmb{i}}^{\bar{x}}\lambda_{\pmb{i}}^{y},\qquad
\sigma_{\pmb{i}}^{z}=\mathrm{i}\lambda_{\pmb{i}}^{x}\lambda_{\pmb{i}}^{\bar{x}}
\tag{10.2.5}
\end{equation*}
作用在物理希尔伯特空间之内，且满足泡利矩阵的代数。因此，我们可以把 $\sigma_{\pmb{i}}^{l}$ 认作自旋算符。利用在物理希尔伯特空间之内
\begin{equation*}
(-)^{\psi^{\dagger}_{1,\pmb{i}}\psi_{1,\pmb{i}}+\psi^{\dagger}_{2,\pmb{i}}\psi_{2,\pmb{i}}}
=\lambda_{\pmb{i}}^{x}\lambda_{\pmb{i}}^{y}\lambda_{\pmb{i}}^{\bar{x}}\lambda_{\pmb{i}}^{\bar{y}}=1
\end{equation*}
这一事实，我们可以证明，费米子哈密顿量 (10.2.2) 变为（见问题 10.2.2）
\begin{equation*}
H_{\text{spin}}=-\sum_{\pmb{i}}V_{\pmb{i}}\hat F_{\pmb{i}},\qquad
\hat F_{\pmb{i}}=\sigma_{\pmb{i}}^{x}\sigma_{\pmb{i}+\hat{\pmb{x}}}^{y}\sigma_{\pmb{i}+\hat{\pmb{x}}+\hat{\pmb{y}}}^{x}\sigma_{\pmb{i}+\hat{\pmb{y}}}^{y}
\tag{10.2.6}
\end{equation*}
在物理希尔伯特空间之内。

物理希尔伯特空间中的所有态（即自旋 $1/2$ 模型中的所有态）都可以从诸 $|\{s_{ij}\}\rangle$ 态投影到物理希尔伯特空间而得到，即 $\mathcal{P}|\{s_{ij}\}\rangle$。投影算符为
\begin{equation*}
\mathcal{P}=\prod_{\pmb{i}}\frac{1+(-)^{\psi^{\dagger}_{1\pmb{i}}\psi_{1\pmb{i}}+\psi^{\dagger}_{2\pmb{i}}\psi_{2\pmb{i}}}}{2}
\end{equation*}
由于 $[\mathcal{P},H]=0$，投影后的态 $\mathcal{P}|\{s_{ij}\}\rangle$（如果非零）仍然是 $H$（或 $H_{\text{spin}}$）的本征态，且具有相同的本征值。经过投影，相互作用费米子模型的精确解导致自旋 $1/2$ 模型的精确解。

我们注意到，对于 $x$ 和 $y$ 两个方向都有周期性边界条件的系统，所有链的乘积
\begin{equation*}
\prod_{\pmb{i}}\,(\mathrm{i}s_{\pmb{i},\pmb{i}+\pmb{x}})(\mathrm{i}s_{\pmb{i},\pmb{i}+\pmb{y}})=(-)^{\hat N_f},
\tag{10.2.7}
\end{equation*}
其中 $\hat N_f=\sum_{\pmb{i}}(\psi^{\dagger}_{1\pmb{i}}\psi_{1\pmb{i}}+\psi^{\dagger}_{2\pmb{i}}\psi_{2\pmb{i}})$ 是总费米子数算符。因此，$|\{s_{ij}\}\rangle$ 的投影非零，仅当
\begin{equation*}
\prod_{\pmb{i}}\,(\mathrm{i}s_{\pmb{i},\pmb{i}+\pmb{x}})(\mathrm{i}s_{\pmb{i},\pmb{i}+\pmb{y}})=1.
\tag{10.2.8}
\end{equation*}

物理态（每个格点具有偶数个费米子）在由
\begin{equation*}
\hat G=\prod_{\pmb{i}}G_{\pmb{i}}^{\;\psi^{\dagger}_{1,\pmb{i}}\psi_{1,\pmb{i}}+\psi^{\dagger}_{2,\pmb{i}}\psi_{2,\pmb{i}}}
\end{equation*}
生成的局域 $Z_2$ 变换下不变，其中 $G_{\pmb{i}}$ 是一个只取 $\pm1$ 两个值的任意函数。我们注意到，$Z_2$ 变换把 $\psi_{I\pmb{i}}$ 变为 $\tilde\psi_{I\pmb{i}}=G_{\pmb{i}}\psi_{I\pmb{i}}$，把 $s_{ij}$ 变为 $\tilde s_{ij}=G_{\pmb{i}}s_{ij}G_{\pmb{j}}$。我们发现，$|\{s_{ij}\}\rangle$ 与 $|\{\tilde s_{ij}\}\rangle$ 在投影后给出同一个物理态（如果它们的投影非零）。于是，$\{s_{ij}\}$ 是自旋 $1/2$ 系统 $H_{\text{spin}}$ 的精确本征态的一个多对一标签。向具有每格点偶数个费米子的物理希尔伯特空间的投影，使我们的理论成为一个 $Z_2$ 规范理论。

\subsection*{问题 10.2.1}

\begin{enumerate}
\item[(a)] 证明式 (10.2.1) 中的诸 $\hat U_{ij}$ 构成一组彼此对易的算符。
\item[(b)] 证明式 (10.2.3) 中的诸 $\psi$ 满足复费米子的标准反对易关系。
\end{enumerate}

\subsection*{问题 10.2.2}

\begin{enumerate}
\item[(a)] 证明 $2\psi_1^\dagger\psi_1-1=\mathrm{i}\lambda^{x}\lambda^{\bar{x}}$ 以及 $2\psi_2^\dagger\psi_2-1=\mathrm{i}\lambda^{y}\lambda^{\bar{y}}$。证明在一个格点上具有偶数个费米子的态是 $\lambda^{x}\lambda^{y}\lambda^{\bar{x}}\lambda^{\bar{y}}$ 的本征值为 $+1$ 的本征态。
\item[(b)] 证明式 (10.2.7)。
\item[(c)] 证明式 (10.2.5) 中的诸 $\sigma^{l}$ 在物理希尔伯特空间内满足泡利矩阵的代数。
\item[(d)] 证明费米子哈密顿量 (10.2.2) 在物理希尔伯特空间内变为自旋哈密顿量 (10.2.6)。
\end{enumerate}

\subsection*{问题 10.2.3}

我们已经看到，把费米子哈密顿量 $H$ 限制在每格点具有偶数个费米子的子空间之内，可以从费米子系统得到一个自旋 $1/2$ 系统。我们注意到，费米子哈密顿量 $H$ 也作用在每格点具有奇数个费米子的子空间之内。试把费米子哈密顿量 $H$ 限制在奇数费米子子空间之内，得到相应的自旋 $1/2$ 系统。

\subsection{精确本征态与拓扑简并的基态}

\begin{keypoints}
\item 基态简并受到 $Z_2$ 规范结构的保护，对任何局域微扰都是稳健的。
\end{keypoints}

上述 $Z_2$ 规范结构使我们能够清点由诸 $|\{s_{ij}\}\rangle$ 态投影得到的物理态的数目。我们再次假定两个方向都有周期性边界条件。注意到常数 $Z_2$ 规范变换 $G_{\pmb{i}}=-1$ 不改变 $s_{ij}$，我们看到，彼此规范等价的互不相同的 $s_{ij}$ 只有 $2^{N_{\text{site}}}/2$ 个。在 $4^{N_{\text{site}}}$ 个 $|\{s_{ij}\}\rangle$ 态中，有 $4^{N_{\text{site}}}/2$ 个满足 $\prod_{\pmb{i}}(\mathrm{i}s_{\pmb{i},\pmb{i}+\pmb{x}})(\mathrm{i}s_{\pmb{i},\pmb{i}+\pmb{y}})=1$（即具有偶数个费米子）。因此，诸 $|\{s_{ij}\}\rangle$ 态的投影至多给我们 $(4^{N_{\text{site}}}/2)\allowbreak/(2^{N_{\text{site}}}/2)\allowbreak=2^{N_{\text{site}}}$ 个物理态。另一方面，所有 $|\{s_{ij}\}\rangle$ 态的投影应当给我们全部 $2^{N_{\text{site}}}$ 个自旋态。因此，不同的 $s_{ij}$ 的 $Z_2$ 规范等价类必须导致独立的物理态，这样投影才能恢复全部 $2^{N_{\text{site}}}$ 个自旋态。$s_{ij}$ 的 $Z_2$ 规范等价类是周期格子上所有物理态的一对一标签。这样，我们能够得到周期格子上 $H_{\text{spin}}$ 的所有本征态和本征值。

我们注意到，$Z_2$ 通量 $F_{\pmb{i}}$ 在 $Z_2$ 规范变换下不变。因此，如果两个位形 $s_{ij}$ 与 $\tilde s_{ij}$ 分别具有不同的 $Z_2$ 通量 $F_{\pmb{i}}$ 与 $\tilde F_{\pmb{i}}$，则 $s_{ij}$ 与 $\tilde s_{ij}$ 必定属于两个不同的 $Z_2$ 规范等价类。两个态 $|s_{ij}\rangle$ 与 $|\tilde s_{ij}\rangle$ 在投影后将导致不同的物理态。这一结果提示，投影后的态（物理态）用 $Z_2$ 通量 $F_{\pmb{i}}$ 来标记更好。结果发现，$F_{\pmb{i}}$ 并不是物理态的一个完美的一对一标签。对于周期性边界条件下的偶 $\times$ 偶格子上的系统，每个 $F_{\pmb{i}}$ 标记 4 个态（即对每一个 $Z_2$ 通量位形 $F_{\pmb{i}}$，有且仅有 4 个物理态给出同一个 $F_{\pmb{i}}$）。具有相同 $F_{\pmb{i}}$ 的 4 个物理态全部具有相同的能量。

为了理解这一四重简并，让我们考虑由某个 $|\{s_{ij}\}\rangle$ 态的投影所给出的一个本征态。其他简并的本征态可以通过施行下列两个变换得到：
\begin{align*}
T_1&:(s_{\pmb{i},\pmb{i}+\pmb{x}},s_{\pmb{i},\pmb{i}+\pmb{y}})\to(s_{\pmb{i},\pmb{i}+\pmb{x}},(-)^{\delta_{iy}}s_{\pmb{i},\pmb{i}+\pmb{y}})\\
T_2&:(s_{\pmb{i},\pmb{i}+\pmb{x}},s_{\pmb{i},\pmb{i}+\pmb{y}})\to((-)^{\delta_{ix}}s_{\pmb{i},\pmb{i}+\pmb{x}},s_{\pmb{i},\pmb{i}+\pmb{y}})
\end{align*}
我们看到，$T_1$ 不改变 $s_{\pmb{i},\pmb{i}+\pmb{x}}$，只在 $i_y=0$ 时翻转 $s_{\pmb{i},\pmb{i}+\pmb{y}}$ 的符号。由构造可知，$T_1$ 与 $T_2$ 不改变 $Z_2$ 通量 $F_{\pmb{i}}$。如果我们把周期格子看作一个环面，那么 $T_1$ 与 $T_2$ 通过环面的两个洞各插入 $\pi$ 通量（见图 9.6）。

在偶 $\times$ 偶格子上，变换 $T_1$ 与 $T_2$ 不改变乘积 $\prod_{\pmb{i}}(\mathrm{i}s_{\pmb{i},\pmb{i}+\pmb{x}})(\mathrm{i}s_{\pmb{i},\pmb{i}+\pmb{y}})$。因此，三个变换 $T_1$、$T_2$ 与 $T_1T_2$ 生成另外三个简并态。我们看到，所有能量本征值都具有四重简并。特别地，自旋 $1/2$ 模型 $H_{\text{spin}}$ 在偶 $\times$ 偶周期格子上具有 4 个简并基态。

在偶 $\times$ 奇格子上，情况稍有不同。$T_2$ 生成的态具有奇数个费米子，不对应于任何物理自旋 $1/2$ 态。因此，我们只能用 $T_1$ 来生成另一个简并态。在偶 $\times$ 奇周期格子上只有两个简并基态（由 $T_1$ 生成）。在奇 $\times$ 奇格子上，也有由 $T_1T_2$ 生成的两个简并基态。

我们注意到，局域地看，$T_1$ 与 $T_2$ 变换与 $Z_2$ 规范变换不可区分。由于物理自旋算符在 $Z_2$ 规范变换下不变，它们在 $T_1$ 与 $T_2$ 变换下也不变。因此，即使我们在精确可解模型 (10.2.6) 上加入任意局域微扰，由 $T_1$ 与 $T_2$ 生成的简并基态仍将保持简并。基态的这一简并是一个稳健的拓扑性质，表明基态中存在非平凡的拓扑序。

\subsection*{问题 10.2.4}

考虑 $L_x\times L_y$ 周期格子上的式 (10.2.6)。我们假定 $V_{\pmb{i}}=V$，只在某一行方格上 $V_{\pmb{i}}=0$。这时可以把该系统看作定义在一个圆柱上，圆柱在 $x$ 方向具有两条圆形边缘。试求系统的基态简并度。这些简并基态可以看作边缘态。证明每个边缘格点有 $\sqrt{2}$ 个边缘态。这意味着边缘激发由一个马约拉纳费米子描写。

\subsection{基态的投影对称群表征}

\begin{keypoints}
\item 式 (10.2.6) 的 $V<0$ 与 $V>0$ 基态具有不同的量子序。
\end{keypoints}

在本节中，我们将假定自旋 $1/2$ 系统 (10.2.6) 中的 $V_{\pmb{i}}$ 是均匀的，即 $V_{\pmb{i}}=V$。我们已经通过把自旋算符写成费米子算符的乘积求解了式 (10.2.6)（见式 (10.2.5)），也就是说，我们用投影构造求解了自旋 $1/2$ 系统。十分有趣的是，对于这个特定的自旋 $1/2$ 系统 (10.2.6)，投影构造竟然给出精确的结果。由于精确基态波函数是由投影构造得到的（即通过投影自由费米子波函数 (9.2.11) 而得到），我们可以用第 9 章发展的 PSG 表征方法来表征基态中的量子序。

第 9 章的讨论限于自旋旋转不变的态。下面，我们将把平均场理论和 PSG 推广到自旋旋转非不变的态。如果我们把 $|\downarrow\rangle$ 态认作零玻色子态 $|0\rangle$，把 $|\uparrow\rangle$ 态认作单玻色子态 $|1\rangle$，那么自旋 $1/2$ 模型 $H_{\text{spin}}$ 也可以看作一个硬核玻色子模型。下面我们将用玻色子图像来描写我们的模型。

为了用投影构造构造量子有序（或纠缠）的多玻色子波函数，我们首先引入“平均场”费米子哈密顿量
\begin{equation*}
H_{\text{mean}}=\sum_{\langle ij\rangle}\left(\psi^{\dagger}_{I,\pmb{i}}u_{ij}^{IJ}\psi_{J,\pmb{j}}+\psi^{\dagger}_{I,\pmb{i}}w_{ij}^{IJ}\psi^{\dagger}_{J,\pmb{j}}+\text{h.c.}\right)
\tag{10.2.9}
\end{equation*}
其中 $I,J=1,2$。我们将用 $u_{ij}$ 和 $w_{ij}$ 表示元素分别为 $u_{ij}^{IJ}$ 和 $w_{ij}^{IJ}$ 的 $2\times2$ 复矩阵。令 $|\Psi^{(u_{ij},w_{ij})}_{\text{mean}}\rangle$ 为上述自由费米子哈密顿量的基态，则下列多体玻色子波函数可以得到：
\begin{equation*}
\Phi^{(u_{ij},w_{ij})}(\pmb{i}_1,\pmb{i}_2...)=\langle0|\prod_{n}b(\pmb{i}_n)|\Psi^{(u_{ij},w_{ij})}_{\text{mean}}\rangle
\tag{10.2.10}
\end{equation*}
其中
\begin{equation*}
b(\pmb{i})=\psi_{1,\pmb{i}}\psi_{2,\pmb{i}}
\tag{10.2.11}
\end{equation*}
我们注意到，物理玻色子波函数 $\Phi^{(u_{ij},w_{ij})}(\{\pmb{i}_n\})$ 在下列 $SU(2)$ 规范变换下不变：
\begin{equation*}
(\psi_{\pmb{i}},u_{ij},w_{ij})\to(G(\pmb{i})\psi_{\pmb{i}},\,G(\pmb{i})u_{ij}G^{\dagger}(\pmb{j}),\,G(\pmb{i})w_{ij}G^{\top}(\pmb{j}))
\end{equation*}
其中 $G(\pmb{i})\in SU(2)$。

玻色子波函数 $\Phi^{(u_{ij},w_{ij})}(\{\pmb{i}_n\})$ 中的量子序可以用 PSG 来表征。PSG 由使拟设 $(u_{ij},w_{ij})$ 不变的对称变换与 $SU(2)$ 规范变换的联合变换构成。

这里我们想指出，诸 $\hat U_{ij}$ 的共同本征态，即 $|\{s_{ij}\}\rangle$，是自由费米子系统
\begin{equation*}
\tilde H_{\text{mean}}=\sum_{\langle ij\rangle}\left(s_{ij}\hat U_{ij}+\text{h.c.}\right)
\tag{10.2.12}
\end{equation*}
的基态。因此，从投影构造的观点看，$\tilde H_{\text{mean}}$ 可以看作“平均场”哈密顿量 $H_{\text{mean}}$。事实上，$\tilde H_{\text{mean}}$ 是 $H_{\text{mean}}$ 的一个特例，其中
\begin{equation*}
-w_{i,i+\hat{\pmb{x}}}=u_{i,i+\hat{\pmb{x}}}=-\mathrm{i}s_{i,i+\hat{\pmb{x}}}(1+\tau^{3}),\qquad
-w_{i,i+\hat{\pmb{y}}}=u_{i,i+\hat{\pmb{y}}}=-\mathrm{i}s_{i,i+\hat{\pmb{y}}}(1-\tau^{3})
\end{equation*}
如果 $s_{ij}$ 与 $(u_{ij},w_{ij})$ 通过上式相联系，则态 $|\{s_{ij}\}\rangle$ 就等于平均场态 $|\Psi^{(u_{ij},w_{ij})}_{\text{mean}}\rangle$。物理自旋波函数由平均场态的投影 $\mathcal{P}|\{s_{ij}\}\rangle$ 得到，这与式 (10.2.10) 等价。

记住，在投影构造中，我们需要选择拟设 $u_{ij}$ 与 $w_{ij}$，使平均能量 $\langle\Psi^{(u_{ij},w_{ij})}_{\text{mean}}|\mathcal{P}H_{\text{spin}}\mathcal{P}|\Psi^{(u_{ij},w_{ij})}_{\text{mean}}\rangle$ 极小。对我们的自旋模型 $H_{\text{spin}}$，这样得到的平庸基态恰好就是精确基态！而且，如果我们选择一个不同的 $s_{ij}$，投影态 $\mathcal{P}|\Psi^{(u_{ij},w_{ij})}_{\text{mean}}\rangle$ 将对应于 $H_{\text{spin}}$ 的一个具有能量 (10.2.4) 的精确本征态。正是在这个意义上，投影构造提供了 $H_{\text{spin}}$ 的一个精确解。

当 $V>0$ 时，我们模型的基态由 $Z_2$ 通量位形 $F_{\pmb{i}}=1$ 给出。为了产生这样的通量，我们可以选择 $s_{\pmb{i},\pmb{i}+\pmb{x}}=s_{\pmb{i},\pmb{i}+\pmb{y}}=\mathrm{i}$。这时，式 (10.2.12) 变为带有
\begin{equation*}
-w_{i,i+\pmb{x}}=u_{i,i+\pmb{x}}=1+\tau^{3},\qquad
-w_{i,i+\pmb{y}}=u_{i,i+\pmb{y}}=1-\tau^{3}.
\end{equation*}
的式 (10.2.9)。为得到上述拟设的 IGG，我们注意到 $u_{ij}$ 在常数规范变换 $G(\pmb{i})=\mathrm{e}^{\mathrm{i}\phi\tau^{z}}$ 下不变。然而，$w_{ij}$ 仅当 $\phi=0,\pi$ 时才不变。于是，IGG $=Z_2$。我们发现低能有效理论是一个 $Z_2$ 规范理论。由于该拟设已经是平移不变的，它在带平凡规范变换 $G_x=G_y=\tau^{0}$ 的 $G_xT_x$ 与 $G_yT_y$ 下不变。上述拟设的 PSG 就是式 (9.4.10) 中的 Z2A PSG。因此，$V>0$ 的基态是一个 Z2A 态。注意，对 Z2A PSG，$G_xT_x$ 与 $G_yT_y$ 满足平移代数 $(G_xT_x)^{-1}(G_yT_y)^{-1}(G_xT_x)(G_yT_y)=1$。

当 $V<0$ 时，基态由位形 $F_{\pmb{i}}=-1$ 给出，它可以由 $s_{\pmb{i},\pmb{i}+\pmb{x}}=(-)^{i_y}s_{\pmb{i},\pmb{i}+\pmb{y}}=\mathrm{i}$ 产生（译注：按本页下方拟设中 $y$ 方向链上的 $(-)^{i_x}$ 因子以及 $F_{\pmb{i}}=-1$ 的要求，此处指数应为 $i_x$，即 $s_{\pmb{i},\pmb{i}+\pmb{y}}=(-)^{i_x}\mathrm{i}$；原文印作 $(-)^{i_y}$）。拟设现在具有形式
\begin{equation*}
-w_{i,i+\pmb{x}}=u_{i,i+\pmb{x}}=1+\tau^{3},\qquad
-w_{i,i+\pmb{y}}=u_{i,i+\pmb{y}}=(-)^{i_x}(1-\tau^{3}).
\end{equation*}
该拟设在平移 $\pmb{i}\to\pmb{i}+\pmb{x}$ 接规范变换 $G_x(\pmb{i})=(-)^{i_y}$ 下不变。它的 PSG 是式 (9.4.11) 中的 Z2B PSG。因此，$V<0$ 的基态是一个 Z2B 态。对 Z2B PSG，$G_xT_x$ 与 $G_yT_y$ 满足磁平移代数 $(G_xT_x)^{-1}(G_yT_y)^{-1}(G_xT_x)(G_yT_y)=-1$。费米子跳跃哈密顿量 $H_{\text{mean}}$ 描写费米子在磁场中的跳跃，每个方格有 $\pi$ 通量。不同的 PSG 告诉我们，$V<0$ 与 $V>0$ 的基态具有不同的量子序。

\section{正方格子上的 $Z_2$ 自旋液体与弦网凝聚}

\begin{keypoints}
\item 弦网凝聚导致演生的 $Z_2$ 规范理论和费米子。
\end{keypoints}

在这一节中，我们将从弦网凝聚的观点（Levin and Wen, 2003）出发，讨论精确可解自旋 $1/2$ 模型 (10.2.6)（Kitaev, 2003; Wen, 2003c）。该模型是最简单的模型之一，它演示了局域玻色模型中弦网凝聚与演生规范玻色子及费米子之间的联系。该模型也可以用从属玻色子方法求解，这使我们能够看到描写量子序的 PSG 如何与弦网凝聚相联系（见 10.4 节）。

\subsection{构造具有闭合弦网凝聚的哈密顿量}

\begin{keypoints}
\item 可以通过让某些弦状物体凝聚来构造具有演生规范场的局域玻色模型。
\end{keypoints}

让我们首先考虑正方格子上一个任意的自旋 $1/2$ 模型。我们想问的第一个问题是：什么样的自旋相互作用能够给出一个低能规范理论？如果我们相信规范理论与弦网理论之间的联系（Banks et al., 1977; Savit, 1980; Wen, 2003a），那么得到低能规范理论的一种途径，是设计一种自旋相互作用，使它允许大闭合弦的强烈涨落，却禁止其他类型的涨落（例如局域自旋翻转、开弦涨落等）。我们希望，大闭合弦的强烈涨落的出现将导致任意大小闭合弦的凝聚，继而给出一个低能规范理论。

%FIG{10.1}: {$H_J$ 的基态之上的一个开弦激发。}

让我们从
\begin{equation*}
H_J=-J\sum_{\text{even}}\sigma_{\pmb{i}}^{x}-J\sum_{\text{odd}}\sigma_{\pmb{i}}^{y}
\end{equation*}
出发，其中 $\pmb{i}=(i_x,i_y)$ 标记格点，$\sigma^{x,y,z}$ 是泡利矩阵，而 $\sum_{\text{even}}$（或 $\sum_{\text{odd}}$）对偶格点求和，偶格点满足 $(-)^{\pmb{i}}\equiv(-1)^{i_x+i_y}=1$（或对奇格点求和，奇格点满足 $(-)^{\pmb{i}}=-1$）。$H_J$ 的基态，即 $|0\rangle$，具有这样的性质：偶格点上的自旋指向 $x$ 方向，奇格点上的自旋指向 $y$ 方向（见图 10.1）。这样一个态将被定义为无弦态（no-string state）。

为了产生一个弦激发，我们先画一条连接最近邻偶方格的弦（见图 10.1），然后把弦上的自旋翻转。这样一个弦态由下列弦产生算符（string creation operator，或简称弦算符）产生：
\begin{equation*}
W(C)=-\prod_{C}\sigma_{\pmb{i}}^{a_{\pmb{i}}}
\tag{10.3.1}
\end{equation*}
其中乘积 $\prod_{C}$ 遍及弦上的所有格点，且若 $\pmb{i}$ 为偶则 $a_{\pmb{i}}=y$，若 $\pmb{i}$ 为奇则 $a_{\pmb{i}}=x$。一个一般的弦态具有形式
\begin{equation*}
|C_1C_2...\rangle=W(C_1)W(C_2)...|0\rangle
\end{equation*}
其中 $C_1,C_2,\dots$ 是弦。这样的态将被称为弦网态（string-net state），因为弦可以相交和重叠。产生弦网的算符，即
\begin{equation*}
W(C_{\text{net}})=W(C_1)W(C_2)...
\end{equation*}
% （本块止于书页 452 末行显示公式；其后半句“将被称为弦网算符。”在书页 453 (p468) 顶部，由下一块收尾。）
